\documentclass[11pt,a4paper]{article}
\usepackage[T1]{fontenc}
\usepackage{lmodern,amsmath,amssymb,amsthm,mathtools,microtype}
\usepackage[margin=27mm]{geometry}
\usepackage[colorlinks=true,linkcolor=blue,urlcolor=blue]{hyperref}
\newtheorem{proposition}{Proposition}
\newtheorem{lemma}[proposition]{Lemma}
\newcommand{\Q}{\mathbb Q}
\title{Two rational sections of the cubic trace family}
\author{}\date{17 November 2026}
\begin{document}\maketitle

Put
\[
 C(V,w)=V^2+w(2-V),\qquad
 E_w:\ y^2=x^3-(2+24w)x^2+(1+1176w)x.
\]
The substitutions $x=1+24V$ and $y=24Y$ identify $E_w$ with
$Y^2=(1+24V)C(V,w)$. It has the rational reference section
$P=(49,336)$, corresponding to $(V,Y)=(2,14)$.

\begin{proposition}\label{prop:sections}
Over $\Q(r)$, put $w=r^2-1$. The points
\[
                 P=(49,336),\qquad Q=(25,120r)
\]
are independent non-torsion points on $E_{r^2-1}$. The point $Q$
corresponds to $(V,Y)=(1,5r)$. Their sums satisfy
\[
 x(P+Q)=(7r-10)^2,\qquad x(P-Q)=(7r+10)^2.
\]
In particular
\[
 V(P+Q)=\frac{(7r-11)(7r-9)}{24}.
\]
\end{proposition}
\begin{proof}
Substitution verifies both points. The slopes for addition are
$14-5r$ and $14+5r$, respectively; the formulas follow from
the Weierstrass group law.

At $r=2$, the curve is
\[
                 y^2=x^3-74x^2+3529x.
\]
It has good reduction at $7$, $11$ and $13$, with respectively
$12$, $16$ and $18$ rational points over those residue fields.
Consequently its rational torsion subgroup has order dividing two.
Both specialized points have nonzero ordinate, so both have infinite
order. Therefore the original sections are non-torsion.

The involution $r\mapsto-r$ fixes $P$ and sends $Q$ to $-Q$.
If $mP+nQ=0$ over $\Q(r)$, applying the involution and subtracting
gives $2nQ=0$, hence $n=0$ and then $m=0$.
\end{proof}

The point-count argument uses the injectivity of rational torsion
under good reduction at an odd prime. Equivalently, one may use
prime-to-residue-characteristic injectivity at the three distinct
primes: it gives the same divisor two for every primary torsion part.
The exact point counts and addition identities are certified by
\texttt{check\_second\_trace\_section.py}.

\section*{Physical normalization and local preparation}

For a rational weight parameter $u\ne1$, define
\[
             w=\frac{4u}{(1-u)^2},\qquad
             r=\frac{1+u}{1-u}.
\]
Then $1+w=r^2$. At an odd prime, if $v_p(u)\ne0$,
\[
                         v_p(w)=|v_p(u)|;
\]
if $v_p(u)=0$, the valuation of $w$ is even. Thus every odd
valuation of $u$ becomes a positive odd valuation of $w$.
This observation does not on its own transfer a full detector:
its scalar factors, height bounds, and the denominator $1-u$
must all be treated coherently.

As $u$ tends to zero at a fixed place, $w$ tends to zero and
$r$ tends to one. Both sections stay rational. The reference section
has trace two and the second section has trace one, exactly the two
residues permitted by the reciprocal detector tests at selected
contacts. At a hypothetical bad pole, the multiplier $1+24V$
still reduces to one. Consequently adjoining the second section
does not change the local square-class argument.

The sections do not identify an actual detector trace. In particular,
forcing an actual trace to equal the constant value one is a separate
and generally obstructed condition. Likewise, generic independence
of two sections does not prove independent specialization at every
rational $r$, or simultaneous approximation to independently chosen
points at all places.

\section*{The remaining incidence condition}

For actual detector points and a physical parameter $u$, write
$V=\operatorname{Tr}(W^{-1})$ and $r=(1+u)/(1-u)$.
A sufficient completeness construction would produce an integer
pair $(a,b)$ such that
\[
                    1+24V=x(aP+bQ)
\]
on $E_{r^2-1}$, together with every original point, difference,
height and nonzero constraint. The existence of the sections proves
that the right-hand side has many rational values. It does not
prove this prescribed-abscissa equality. No reduction to twenty-one
or nine defining coordinates follows from this proposition alone.
\end{document}
